\subsection{Single-node Ibex GEMM evaluation}
\label{sec:ibex-gemm}
The GEMM campaign compares RCOMPSs with PaRSEC/DPLASMA on one allocated Ibex node containing four matching NVIDIA A100 GPUs. Configurations use one, two, or four of those GPUs with the same allocation of schedulable physical CPU cores. Scheduler-reserved cores are excluded, and BLAS calls are single-threaded to avoid nested oversubscription. The recorded node model, GPU capacities, CPU mask, driver, CUDA libraries, source revisions, and compatibility patches define the final testbed. Different nodes or GPU models will not be pooled into one scaling curve.

\paragraph{Work and timing contract.}
Both systems evaluate square FP64 GEMM with $\alpha=0.51$ and $\beta=-0.42$. The primary comparison measures fresh-process completion, including runtime initialization, input generation, task execution, transfers, complete host output, and shutdown. RCOMPSs returns every output block and assembles the complete R matrix on the master. Native-operation and PaRSEC enqueue/progress/destruction timers are supplementary diagnostics with explicitly different boundaries. Tile layout and input generators differ between implementations; the evaluation compares complete implementations rather than attributing every difference to the runtime alone.

RCOMPSs generates reproducible block inputs from global matrix indices and assigns disjoint output regions to CPU and GPU tasks. An independent small calibration determines their static work fractions; the fraction is fixed before each reported case and is included in the artifact. This application-level partition is distinct from COMPSs resource scheduling. The benchmark explicitly selects the allocated CUDA ordinal before creating a context in a persistent executor, and records the observed PCI bus and CPU affinity. Repeated concurrent preflight tasks must demonstrate all requested GPUs before multi-GPU results are accepted.

\paragraph{Size, capacity, and CPU studies.}
The size sweep includes small matrices, larger compute-intensive cases, and points near the measured single-device allocation boundary. The artifact distinguishes three full resident matrices from a streamed reduction with resident $C$ and from bounded output-block tasks. Actual allocation probes reserve headroom and report the largest tested success; they do not establish GEMM correctness at that dimension or an absolute allocator limit. A separately labeled resident GEMM at a conservatively rounded fitting size checks sampled dot products and records failures; this one-run capacity demonstration is not a repeated performance comparison. Thus a streamed calculation larger than the resident-data boundary is not described as fitting all inputs on one GPU. A separate CPU-core sweep uses a fixed matrix dimension, with CPU-only and one-A100 panels. This distinguishes CPU task parallelism from host-thread support for GPU execution.

\paragraph{Validation and memory.}
Each RCOMPSs output block is checked using independent sampled dot products with higher-precision accumulation, including partial edge blocks; small cases receive full-output checks. Large sampled checks are not described as full-matrix validation. PaRSEC receives numerical preflight checks, while production-size output validation is reported separately. Every timing point has one excluded warmup and three measured repetitions with alternating runtime order; min--max whiskers describe observed variability, not confidence intervals.

Peak host memory is the simultaneous sum of proportional set size (PSS) over case-tagged master and worker processes, including workers launched through SSH. This avoids double-counting shared pages as summed RSS would. GPU memory is sampled per process and device; the aggregate curve uses the simultaneous device-memory sum rather than adding independently observed device peaks. Both include runtime contexts, workspaces, and any device-cache reservations. Lower observed allocation alone is not proof of lower live-data complexity or fewer transfers; caching policy can trade memory for performance. Sampling every 0.5~s provides an observed peak, not a guarantee of capturing arbitrarily short allocations. Allocation probes additionally record CUDA allocator measurements directly. Failed, incomplete, out-of-memory, or incorrectly bound cases are labeled and excluded from successful timing curves.

\begin{figure*}[t]\centering
\IfFileExists{Figures/ibex_gemm_scaling.pdf}{\includegraphics[width=0.94\textwidth]{Figures/ibex_gemm_scaling.pdf}}{\plotbox{Ibex GEMM size and GPU scaling --- pending}{One node; RCOMPSs and PaRSEC; 1, 2, and 4 matching A100 GPUs; identical physical CPU allocation. X: matrix dimension. Y: fresh-process completion time. A capacity marker identifies the separately tested three-matrix allocation point.}{Figures/ibex_gemm_scaling.pdf}}
\caption{Single-node size scaling. Measured curves will be inserted after numerical, device-assignment, and completeness checks pass. The allocation marker is not a compute limit for streamed tasks.}\label{fig:ibex-gemm-scaling}\end{figure*}

\begin{figure*}[t]\centering
\IfFileExists{Figures/ibex_gemm_cpu_threads.pdf}{\includegraphics[width=0.94\textwidth]{Figures/ibex_gemm_cpu_threads.pdf}}{\plotbox{Ibex CPU-core sensitivity --- pending}{Fixed matrix dimension. Left: CPU-only task execution. Right: one A100 with varying physical CPU cores. Both systems use the same core masks and single-threaded BLAS calls.}{Figures/ibex_gemm_cpu_threads.pdf}}
\caption{CPU-core sensitivity for CPU-only and one-GPU execution on the same node. Allocated cores are distinguished from instantaneous utilization.}\label{fig:ibex-gemm-cpu}\end{figure*}

\begin{figure*}[t]\centering
\IfFileExists{Figures/ibex_gemm_memory.pdf}{\includegraphics[width=0.94\textwidth]{Figures/ibex_gemm_memory.pdf}}{\plotbox{Ibex host and GPU memory --- pending}{X: matrix dimension. Left: peak summed host PSS. Right: peak simultaneous device-memory sum. Curves distinguish RCOMPSs/PaRSEC and 1/2/4 GPUs.}{Figures/ibex_gemm_memory.pdf}}
\caption{Observed memory consumption for the size-scaling cases. Host and device memory are separate quantities; process contexts, workspaces, and sampling resolution are included in the interpretation.}\label{fig:ibex-gemm-memory}\end{figure*}
